\frontmatter
\begin{titlepage}
  \centering
  \vspace*{0.18\textheight}
  {\zihao{1}\bfseries 高中物理竞赛一轮讲义\par}
  \vspace{1.2em}
  {\zihao{2}\color{DeepBlue}\bfseries 力学补编·运动学\par}
  \vspace{2.6em}
  {\large 依据原《资料.pdf》第二章完整整理\par}
  \vfill
  {\small\color{TextGray}完整正文、109道例题解答与全部重绘教学图\par}
\end{titlepage}

\tableofcontents
\clearpage

\chapter*{序言——写在前面}
\addcontentsline{toc}{chapter}{序言——写在前面}

\section*{什么是物理学？}

相信大家经过小学自然科学和初中物理的学习，对物理已经有了或多或少的了解。但是，在进入高中物理学习，尤其是进入高中物理竞赛学习之前，有必要认真思考“物理学”这个概念本身。这有助于确定研究范围、选择研究方法，也有助于理解公式背后的思想。

物理首先是一门自然科学（nature science），它研究自然现象；这一点把物理学同关注人和社会的人文学科区分开来，也把它同纯粹抽象的数学区分开来。再看研究方法：物理学起源于古希腊的自然哲学（nature philosophy），是一门依赖经验的科学，大胆猜想和定性描述是许多早期物理论述的特征。即使在现代，许多伟大理论的诞生也离不开物理直觉和看似不直观的大胆假设。

\sourceimage{vision-quantum-diptych-v1.png}{图1.1—图1.2（重新创作）：左图表现毕达哥拉斯式“视觉微粒由眼睛射向物体”的历史设想；右图用一颗经典台球分化为多条概率路径，表现海森堡等人讨论的微观粒子弥散。}

另一方面，物理又是实验科学（experimental science）。从近代开始，它越来越强烈地依赖数学的定量逻辑推演和实验验证。尤其在现代物理中，有时并非先有物理猜想再等待数学和实验论证，而是数学结构中的灵感或实验室里的发现直接带动物理理论向前发展。

1928年，Dirac 从相对论量子力学的代数结构出发得到后来以他命名的方程；它一方面解释了电子自旋的若干性质，另一方面预言了电子的反粒子——正电子。电子与正电子相遇时可以湮灭并产生两个高能光子，这一发现翻开了粒子物理的新篇章。再如在2002年和2015年两度获得诺贝尔奖关注的中微子探测实验，确认了中微子振荡，表明中微子质量并非严格为零，从而要求修正人们长期使用的最简标准模型。中微子质量极小、与物质相互作用极弱，因此非常难探测。这些例子说明：物理的发展既可能由实验发现推动，也可能由数学结构引出新的可检验预言。

\PrefaceParticleDiagram

其次要回答：物理包括哪些内容？本讲义采用力学、热学、电磁学、光学和近代物理的顺序，知识由浅入深，并在合适位置引入微积分工具以备后用。这种分类主要服务于学习，不等同于物理学史的时间顺序，也不能代替大学普通物理、更深层次理论或前沿研究中的学科划分。读者在学习讲义的同时，还应尽量了解每一理论的前生今世、适用条件和后续发展，这才有利于进阶学习。

\section*{为什么学习物理？}

最后要回答：为什么学习物理？学习物理意味着用另一种思维方式理解周围世界：把复杂现象抽象成模型，用可测量的量描述，用逻辑和数学推演，再用实验检验。正如英国哲学家 Francis Bacon 所说，读史使人明智，读诗使人聪慧，演算使人精密，哲理使人深刻，道德使人有修养，逻辑修辞使人善辩。物理教育应成为素质教育的一环，使学习者在获得知识的同时提高涵养、形成独立判断。

从长远看，自然科学尤其是物理学的突破会带来科学技术创新。物理既可能成为人类发展的阶梯，也可能成为人类自取灭亡的灾厄。科学史上，Bruno 为坚持信念失去生命，Galileo 身陷囹圄，Nobel 的炸药实验带来家庭悲剧，Marie Curie 长期接触放射性物质并最终罹患疾病。这些故事提醒人们，学术道路有时需要长期付出；但奉献绝不能成为放弃人格、尊严和道德的借口，更不能容许违背学术规范、急功近利、徇私舞弊或放弃诚信与公平。原讲义在这里还借英语谐音留了一句轻松的注脚：“Knowledge is Power, France is Bacon!”——严肃思考之余，也不妨保留一点幽默。

同一理论也可能产生性质完全不同的应用。量子理论催生半导体技术，开启信息时代；核物理既带来能源，也带来毁灭性核武器。因此，学习物理不仅需要知识和能力，也需要诚实、责任和对后果的判断。个人对物理的兴趣可以成为长久动力，但不应以牺牲学术规范、人格与社会责任为代价。

\sourceimage{modern-physics-consequences-v1.png}{图1.4（重新创作）：半导体芯片与核爆场景，表示同一物理理论可能产生截然不同的技术后果。}

希望读者学习物理，不仅因为物理有趣、未知世界吸引着好奇心，也不仅为了享受解决问题的成就感；还应在学习过程中逐渐认清目的，把个人爱好升华为更长远的追求。同时也要全面、均衡发展，培养广泛的思维方式和生活志趣，形成健全人格，做对社会有益的公民。

如果还要问喜欢物理的原因，可以回答：物理世界，异彩纷呈，正像天上的星辰；玄之又玄，众妙之门。理论也许抽象，探索的乐趣却是真实的。

\sourceimage[0.82\linewidth]{physics-board-game-v1.png}{图1.5（重新创作）：原讲义以电视剧《生活大爆炸》中 Sheldon Cooper 设计的“Research Lab”桌游作为趣味情境；此处重新创作为原创的物理研究桌游场景，表示好奇心、交流与游戏化探索。}

\mainmatter
